% ============================================================
% 嵌入式对接教程（二）· Part II：工程实现详解
% 由 emb2_tutorial.tex \input 引入（共用同一份导言区与主题色）
% ============================================================

\part{工程实现详解}

\begin{tipbox}{这一部分写什么}
外设驱动与配置链逐模块讲透：uart\_telemetry 的 bit\_idx 状态机与
v1 重写、uart\_rx\_mon 的"自由运转"监视器、i2c\_bfm\_sc500 的开漏
建模与 ACK 位相遗留、配置链 TB 的对录方法，以及固化 runbook
逐节的执行细节。本册对应迭代记录的第 4、5、6 轮。
\end{tipbox}

% ============================================================
\section{工程总览}

\subsection{外设层文件清单}

\begin{center}
\begin{tabular}{llr l}
\toprule
\rowcolor{rowtint}\textbf{文件} & \textbf{类型} & \textbf{行数} & \textbf{职责} \\
\midrule
\code{rtl/uart\_telemetry.v} & RTL & 94 & 目标状态遥测（115200-8N1） \\
\code{tb/uart\_rx\_mon.v} & 仿真 & 64 & UART 接收监视器（两处复用） \\
\code{tb/tb\_uart\_telemetry.v} & 测试 & 88 & 3 帧注入逐字节验证 \\
\code{rtl/i2c\_bfm\_sc500.v} & 仿真 & 88 & SC500CS I2C 从机 BFM \\
\code{tb/tb\_i2c\_config.v} & 测试 & 93 & 配置链对录（116 条） \\
\midrule
\code{sim/uics500reg\_720p60.v} & 官方 & --- & 116 条寄存器 ROM \\
\code{sim/uii2c.v / uicfgcs500\_720p.v} & 官方 & --- & I2C 控制器与配置引擎 \\
\midrule
\code{scripts/filelist.f} & 脚本 & --- & TD 源文件清单（三段式） \\
\code{scripts/td\_build\_and\_flash.md} & 文档 & --- & 固化 runbook \\
\bottomrule
\end{tabular}
\end{center}

\subsection{外设层的验证策略}

\begin{center}
\begin{tabular}{lll}
\toprule
\rowcolor{rowtint}\textbf{通道} & \textbf{验证方法} & \textbf{状态} \\
\midrule
UART 遥测 & 自收自解（写\code{->}线上\code{->}监视器解码）逐字节比对 & PASS \\
摄像头配置 & 官方配置链 vs BFM 记录，逐事务对录 ROM & 116/116 一致，位相遗留 \\
按键 & 顶层 glue 一行逻辑，随整链仿真覆盖 & 已覆盖 \\
\bottomrule
\end{tabular}
\end{center}

\begin{keybox}{为什么 UART 可以"自收自解"而 I2C 不行}
UART 是\textbf{点对点}：自己发的自己收，闭环即可证明字节正确——
监视器是"独立的第二双眼睛"（独立实现，不共享发送端代码路径）。
I2C 是\textbf{总线}：主控行为（位相/仲裁/时钟拉伸）需要有从机应答
才能走完流程，所以必须写一个 BFM 当对手方——而 BFM 自身的正确性
又需要真机波形佐证（这正是遗留项的由来）。
\end{keybox}

% ============================================================
\section{模块详解：uart\_telemetry}

\subsection{bit\_idx 状态机（v2，重写版）}

发送状态机把"位时间"当节拍：

\begin{Verilog}
if (div_cnt == CLK_DIV - 1) begin          // 一个位时间到
    div_cnt <= 0;
    case (bit_idx)
        4'd0: begin  O_txd <= 1'b0;                 // 起始位
                     shifter <= payload[byte_idx]; end // 装载
        4'd9: begin  O_txd <= 1'b1;                 // 停止位
                     if (byte_idx == NBYTES-1) sending <= 1'b0;
                     else byte_idx <= byte_idx + 1; end
        default: begin O_txd <= shifter[0];          // 数据位 LSB 先行
                       shifter <= {1'b0, shifter[7:1]}; end
    endcase
    bit_idx <= (bit_idx == 4'd9) ? 4'd0 : bit_idx + 4'd1;
end else div_cnt <= div_cnt + 10'd1;
\end{Verilog}

\begin{warnbox}{v1 的失败：装载与移位双写}
第一版把"装载 payload"和"移位输出"写在同一拍的两个分支里，
装载时移位的状态互相覆盖，波形全乱。重写原则：
\textbf{一个节拍只干一件事}——bit\_idx=0 装载（不发数据），
1..8 只移位，9 只出停止位。状态机的每个状态职责单一，
就没有"两件事打架"的空间。
\end{warnbox}

\subsection{组帧快照与跳过}

\begin{Verilog}
tick_d <= I_frame_tick;                        // 边沿检测
if (I_frame_tick && !tick_d && !sending) begin
    payload[0] <= 8'hA5;
    payload[1] <= {7'd0, I_target_valid};
    payload[2] <= I_cx[7:0];   payload[3] <= {5'd0, I_cx[10:8]};
    payload[4] <= I_cy[7:0];   payload[5] <= {5'd0, I_cy[10:8]};
    payload[6] <= {5'd0, I_preset};
    payload[7] <= 8'hA5 + ...（前 6 字节累加和）;
    payload[8] <= 8'h5A;
    sending <= 1'b1;
end
\end{Verilog}

\code{tick\_d} 边沿检测保证"每帧只触发一次"（即使 frame\_tick
因跨域被拉宽）；\code{!sending} 门控实现跳过策略。
校验和在同一拍用组合加法算出来——这拍只有 7 个 8bit 加法，
摆在帧边界上，时序零压力。

\subsection{波特率分频与误差预算}

\begin{Verilog}
CLK_DIV = 646                     // 参数, 74.25e6 / 115200 = 644.5
实际 = 74125000..? 按 74.25e6 / 646 = 114,938 bps
误差 = 114938 / 115200 - 1 = -0.23%   (容差 ±2~3%)
\end{Verilog}

注意这里选的是"整数分频取最接近值"，不追求精确整除——
\textbf{波特率不需要精确，只需要在容差内}。若换 50MHz 时钟，
重算 CLK\_DIV 即可（参数化设计，一处改全局生效）。

% ============================================================
\section{模块详解：uart\_rx\_mon（TB 侧监视器）}

\subsection{两个关键设计}

\begin{Verilog}
// 1) 两级同步器: 串行输入跨进本域
sync <= {sync[0], I_rxd};   wire rxd = sync[1];

// 2) 半位时间后采样, 对齐位中心
if (!rxd && !busy) begin busy <= 1; div <= CLK_DIV/2; ... end
// 3) 假起始重同步: 起始位中心仍为高 -> 判定假起始, 立刻退出
\end{Verilog}

"半位对齐"是 UART 接收的标准姿势：检测到起始\textbf{下降沿}后
等半个位时间再开始按整位采样，采样点就落在每个位的\textbf{中心}，
左右各有半位余量——最抗时钟误差和线上毛刺。

\subsection{v1 的失败：内联监视器漏沿}

第一版在 TB 里内联写监视代码：\code{@(negedge txd)} 等待起始沿、
再阻塞式延迟采数据位。结果字节值逐位错乱——

\begin{warnbox}{事件同步 + 阻塞延迟的经典坑}
\code{@(negedge)} 只对"当前时刻之后"的边沿敏感；而"包内"的
下一位起始沿可能正好发生在阻塞延迟的\textbf{窗口内部}，
被错过之后状态机再也对不齐。修复：重写为\textbf{自由运转状态机}
（每拍看电平、按计数器推进），没有任何 \code{@} 事件同步——
\textbf{时钟节拍是唯一的同步源}。这个教训和 Step3 的
tb\_pwm 重写（任务式同步 -> 连续监视器）是同一课。
\end{warnbox}

\begin{tipbox}{一个监视器，两处复用}
\code{uart\_rx\_mon} 既用在 tb\_uart\_telemetry（验证遥测模块本身），
也用在 tb\_system\_integration（验证整链里的遥测通道）。
独立实现 + 两处复用 = 第三次交叉验证（发送端 FSM、
接收监视器、整链调用者三种视角）。
\end{tipbox}

% ============================================================
\section{模块详解：i2c\_bfm\_sc500}

\subsection{开漏总线的建模}

\begin{Verilog}
assign I_sda = pull ? 1'bz : 1'b0;   // 只能拉低或放手
// TB 侧: pullup(cam_sda);           // "1" 由上拉提供
\end{Verilog}

BFM 不主动驱动高电平——"放手"（z）后由 TB 的上拉（1kΩ 行为模型）
拉高。这是\textbf{忠实建模开漏}的唯一正确方式；
若用 \code{1'b1} 直接驱动，主控和从机同拍驱动相反值时会得到 X，
仿真掩盖真实故障。

\subsection{位相处理与记录}

\begin{Verilog}
wire scl_rise = I_scl && !scl_d;      // 上升沿: 采样/ACK
wire scl_fall = !I_scl && scl_d;      // 下降沿: 数据变化/计数
wire start_c  = I_scl &&  sda_d && !I_sda;   // SCL 高, SDA 下降
wire stop_c   = I_scl && !sda_d &&  I_sda;   // SCL 高, SDA 上升
// 第 9 个上升沿: 地址匹配 -> 拉低 ACK;  第 9 个下降沿: 释放
// 地址匹配: 移位寄存器 == 8'h6C (7bit 0x36 左移 + 写位)
\end{Verilog}

记录数组放在模块内部（\code{reg [7:0] O\_bytes [0:MAX\_BYTES-1]}），
TB 用层次化引用读取——\textbf{Verilog-2001 的端口列表不允许数组}，
这是绕开限制的标准做法（和 system\_top 的遥测层次化引用同源）。

% ============================================================
\section{配置链测试与遗留（诚实的记录）}

\subsection{tb\_i2c\_config 的"对录"方法}

\begin{Verilog}
官方配置引擎 uicfgcs500_720p + uii2c  --驱动--> SCL/SDA
                                              |
                              BFM 记录全部写事务字节
                                              v
期望 ROM uics500reg_720p60 --------------------> 逐条比对
判据: cfg_done 拉高 且 前 116 条写事务 (addr,reg_hi,reg_lo,data)
      与 ROM 完全一致 (AE 追加写只计数不校验)
\end{Verilog}

这个方法的价值：即使 ACK 位相有争议，\textbf{主控侧发出的
寄存器序列本身}已经被独立证明正确（从 ROM 侧逐条确认）——
把"哪一部分可信、哪一部分待定"切得干干净净。

\subsection{遗留：ACK 位相未收敛}

\begin{warnbox}{现象与根因 }
464 次地址识别\textbf{全部正确}（BFM 侧地址匹配逻辑工作正常），
但每字节之后状态被\textbf{自身的 START 检测误判清零}——
BFM 在第 9 个 SCL 上升沿才拉低 SDA 做 ACK，而"SCL 高时 SDA 下降"
正是 START 的定义，于是 ACK 被自己的 START 逻辑吃掉。
两轮位相修复未收敛（改早了波形对不上，改晚了被判 START）。\\[4pt]
\textbf{处置}：明确列为 TD 阶段项——上板用逻辑分析仪抓真机
SCL/SDA 波形对照，半天可收敛；寄存器序列本身已由 ROM 侧确认，
\textbf{外设层的"配置正确性"不依赖这个未收敛项}。
\end{warnbox}

\begin{tipbox}{无硬件条件下的正确姿势}
一个 BFM 位相问题卡住两轮后，最优解不是继续盲试，
而是：\textbf{记录根因假设 + 明确所需证据（真机波形）+ 判定
对交付的影响（不影响）}。工程交付允许带"已知且已界定"的遗留项，
不允许带"说不清"的问题。
\end{tipbox}

% ============================================================
\section{固化 runbook 详解}

\subsection{filelist.f 的三段式}

\begin{Verilog}
# (1) 视觉方向新增: 13 个 rtl (ser96to24 ... vision_top, font_rom)
# (2) 嵌入式新增:  uart_telemetry.v, system_top.v
# (3) 仿真专用 -- 禁止加入 TD 工程:
#     *_stub.v / *_sim.v (替身), sim_video_src.v, frame_dump24.v, tb/*
\end{Verilog}

第（1）段还定义了"哪些视觉文件进综合"：三个生成文件
（\code{hsv\_inv\_lut.v}、\code{font\_rom\_8x16.v}、
\code{marquee\_rom\_64.v}）在内，仿真专用模块不在内。

\subsection{引脚约束增量}

\begin{Verilog}
set_pin_assign {O_pwm_x} -loc {XX} -io_type LVCMOS33
set_pin_assign {O_pwm_y} -loc {XX} -io_type LVCMOS33
set_pin_assign {I_btn_n[0]} -loc {XX} -io_type LVCMOS33
set_pin_assign {I_btn_n[1]} -loc {XX} -io_type LVCMOS33
# 约束文件新增线网: O_pwm_x / O_pwm_y / I_btn_n[1:0] / O_uart_txd
# clock 约束照例程 SDC 不动
\end{Verilog}

脚位从 40Pin DC3 的空闲脚里挑（对照原理图），电平统一 3.3V。
\textbf{舵机信号线接 PWM 脚，舵机电源必须独立 5V 供电并共地}。

\subsection{固化五步与冷启动验收}

\begin{center}
\begin{tabular}{ll}
\toprule
\rowcolor{rowtint}\textbf{步骤} & \textbf{要点} \\
\midrule
1. 打开 BitWizard & 目标器件选\textbf{板载 QSPI Flash}（不是 FPGA SRAM） \\
2. 加载 bit 并转换 & Boot Image Tool，生成 Flash 镜像 \\
3. 烧写 & 用下载器写 Flash \\
4. 断电重上电 & 拔 USB，验证"无 PC 参与" \\
5. 15 秒内出图 & 出图 + 进入跟踪（参数已在 bit 里） \\
\bottomrule
\end{tabular}
\end{center}

\begin{keybox}{五步联调节奏（每步过了才下一步）}
（1）例程原样跑通出图 \code{->}（2）挂 detect\_top：画面无变化 +
ChipWatcher 见 t\_valid \code{->}（3）挂 OSD：出现 TRK0 文字与框
\code{->}（4）接舵机独立电源：duty 变化时画面稳定
\code{->}（5）固化 + 冷启动验收。\\[4pt]
这个顺序的设计逻辑是\textbf{从"官方已知好"逐级加入变量}，
每级只引入一个新东西——出错时责任面永远是一步之宽。
\end{keybox}

% ============================================================
\section{风险清单与交接}

\begin{center}
\begin{tabular}{lp{8.4cm}}
\toprule
\rowcolor{rowtint}\textbf{风险} & \textbf{对策} \\
\midrule
舵机供电 & 独立 5V 共地；画面随舵机动作抖动 = 供电不合格（仿真测不出） \\
AE 拉动颜色 & 演示前锁曝光，或用 8 组阈值现场切换（视觉已备） \\
时序未收敛 & 例程已知现象；挂视觉后重跑时序报告（视觉 <3\% LUT，预期无碍） \\
tuser 链 & 上板"有数据无帧同步"先查 demosaic tuser（第一册界碑） \\
ACK 位相 & 逻辑分析仪抓真机波形，半天收敛（本册遗留项） \\
\bottomrule
\end{tabular}
\end{center}

\begin{okbox}{本册交付清单}
\begin{itemize}
  \item \code{uart\_telemetry.v}：遥测模块（PASS）+ 复用监视器
  \item \code{i2c\_bfm\_sc500.v} + 对录式配置链 TB（116/116 序列确认）
  \item \code{filelist.f} 三段式清单与 \code{td\_build\_and\_flash.md} runbook
  \item 五步联调检查点与四条风险对策（交接给上板阶段）
  \item 一个边界清晰的遗留项（ACK 位相）+ 收敛路径
\end{itemize}
\end{okbox}

% ============================================================
\section{两册总览}

\begin{center}
\begin{tabular}{lll}
\toprule
\rowcolor{rowtint}\textbf{交付层} & \textbf{主张} & \textbf{证据} \\
\midrule
契约层 & 视觉兼容 96bit 契约 & 视觉回归 14 项 \\
挂载层 & 全链可共存（elaboration 零错） & system\_top \\
行为层 & 数据流贯通（受 demosaic 限制） & tb\_system\_integration \\
外设层 & 遥测逐字节正确 & tb\_uart\_telemetry \\
配置层 & 寄存器序列 116 条一致 & tb\_i2c\_config（位相待收敛） \\
固化层 & 上电即运行 & runbook + 五步检查点（上板执行） \\
\bottomrule
\end{tabular}
\end{center}

\begin{tipbox}{嵌入式对接的两册到这里}
第一册管"装进去、跑起来"，第二册管"送出去、配起来、烧进去"。
两册的任务边界就是上板阶段的工作清单——带着这两册和
\code{README\_分工与移交.md} 的勾选表上板，每一步都有据可查。
\end{tipbox}
